\section{The prescribed higher moduli scheme as a finite marking cover}
\label{sec:relative-higher-moduli}

% BEGIN A73H1
\paragraph{A73H1. The geometric fiber at a coprime added level.}
\label{claim:A73H1}
Let $X$ be the supplied lower $\mathrm{Moduli}(D_w,l)$ of RTG1,
$\mathcal A/X$ its universal abelian scheme,
$B=\mathcal O_E$, and $\Lambda=D_w.\mathrm{La}$. Write
$m=l.M$, $\mathfrak b=l.Mb$, $d=\mathrm{tupPrime}(R,t)$, and
$I$ for the prescribed half ideal of RTG10. Thus $m$ is odd,
$m\ge 3$, $(m,d)=1$, $I\bar I=(d)$, and $I+\bar I=B$.
The upper level has modulus $md$, half ideal $\mathfrak b I$
and $A_2=A=1$. Fix an algebraically closed geometric point of $X$,
write $A$ for its fiber and $\eta$ for its lower marking.
The upper markings inducing $\eta$ are in bijection with
\[
 \alpha:\Lambda/I\Lambda\ \xrightarrow{\ \sim\ }_{B/I}\ A[I],
 \qquad
 \eta_2(a)=[e]\eta(a)+\alpha(a\bmod I\Lambda),
 \qquad ed\equiv1\pmod m .
 \tag{A73.1}
\]
Here $A[I]$ denotes the $I$-annihilated summand of $A[d]$,
the group law is written additively, and the auxiliary marking at
$A_2=1$ is zero. For $d=1$ both sides of the proposed isomorphism
are zero and the marking is unique.

\begin{proof}
The lower object's geometric clause gives a single compatible Tate
isomorphism $J$ at \emph{all} integers, preserving its Weil pairings
with its chosen compatible primitive roots. In the splitting of
$\Lambda/md\Lambda$ for the comaximal ideals
$\mathfrak b I$ and $\bar{\mathfrak b}\bar I$, the full $md$-torsion
map of $J$ on the first quotient and zero on the second define a
marking $\eta_2^0$. Indeed on
$\bar{\mathfrak b}\bar I\Lambda$ its value is exactly $J_{md}$,
and it kills $\mathfrak b I\Lambda$ as required.
The $m$-primary part of any upper marking inducing $\eta$ is
$[e]\eta$, since multiplication by $d$ kills the $d$-part and is
invertible on the $m$-part. For $\eta_2^0$ the Tate tower gives
$[d]J_{md}=J_m$, and the upper half idempotent reduces to the
lower one modulo $m$. Hence it induces $\eta$. More explicitly,
if $f m\equiv1\pmod d$ with $d>1$, the two primary components
of $J_{md}(a)$ are $[e]J_m(a)$ and $[f]J_d(a)$ by the tower
identities $[d]J_{md}=J_m$, $[m]J_{md}=J_d$. The same formulas
apply on the marked half after projection. Thus on
$\bar{\mathfrak b}\bar I\Lambda$ the $d$-projection $\alpha_0$
of the default marking satisfies $[m]\alpha_0(a)=J_d(a)$.
For any allowed upper marking, its $d$-primary projection kills
$I\Lambda$ and induces the stated isomorphism onto $A[I]$.
This follows on each primary summand by applying its witnessing
full Tate isomorphism there: the half idempotent is one at the primes
of $I$ and zero at their conjugates. In particular it holds for
$\eta_2^0$, giving a reference $\alpha_0$.

Conversely let $\alpha$ be any of the indicated isomorphisms.
Lift $\alpha_0^{-1}\alpha$ on the $I$-side to an automorphism
of the completed $B$-lattice at each prime dividing $I$. The
completed lattice is free over that discrete valuation ring, and
a matrix lifted from the quotient has unit determinant if its
reduction has unit determinant. Complex conjugation pairs these
primes with disjoint primes of $\bar I$. In the $p$-adic lattice
at a rational $p\mid d$, the integral alternating form pairs the
two factors perfectly: the identities
$B_0(av,w)=B_0(v,\bar a w)$ for $B_0=D_w.EF$ make the off-pair
terms zero and perfection gives duality on the paired summands.
On each conjugate factor take the inverse adjoint of the chosen
automorphism. This preserves the entire alternating form, commutes
with $\mathcal O_E$, and preserves the lattice. Use the identity
on every other factor. The resulting automorphism of the full
profinite lattice induces compatible automorphisms modulo every
positive integer by reduction. At the unused index zero keep
the original additive map. Precomposing $J$ with them keeps its Weil
pairings and roots and its lower marking. Its upper marking has
$d$-projection $\alpha$ and $m$-projection $[e]\eta$, hence is
exactly (A73.1). No lifting of a finite-level matrix to a single
rational unitary matrix is required.
\end{proof}
% END A73H1

% BEGIN A73H2
\paragraph{A73H2. Representability on the actual lower model and rigidity.}
\label{claim:A73H2}
For these prescribed levels there exists an original higher
$\mathrm{Moduli}(D_w,l_2)$ whose forgetful map to the given $X$ is
finite \'etale and surjective. Its underlying scheme is
\[
 P=\operatorname{Isom}_{B/I}
       ((\Lambda/I\Lambda)_X,\mathcal A[I]).
 \tag{A73.2}
\]
The forgetful universal map is cartesian on the abelian schemes,
preserves their rigidified cubical lines, and induces the two
original marking identities of $\mathrm{EtaOK}$.

\begin{proof}
Both levels have the same field
$K_A=\mathbb Q(\mathrm{rt}_0,\zeta_1)$.
The proper semiabelian universal family is abelian, and $d$
is invertible on $X$. Multiplication by $d$ on an abelian scheme
has finite \'etale kernel under this condition. The idempotent of
$B/(d)=(B/I)\times(B/\bar I)$ cuts out its direct summand
$\mathcal A[I]$. Both modules in (A73.2) are finite locally
constant for the \'etale topology. Their module isomorphisms form
a finite locally constant sheaf, hence a finite \'etale scheme:
after trivializing both modules it is the finite set of invertible
linear maps, and the transition permutations give descent.
A73H1 proves nonemptiness at every geometric point. These facts
include the identity cover if $d=1$. Smoothness of the required
relative dimension follows by \'etaleness, and quasiprojectivity
by finiteness over the given quasiprojective scheme. Pull the
universal $\mathrm{SAB}$, its polarization and its rigidification
to $P$ and use its tautological $\alpha$ in (A73.1). Relative
ampleness pulls back, and the polarization kernel stays trivial
by cartesian base change. The determinant
condition persists by base change and A73H1 supplies the full Tate
witness on every geometric fiber. The additivity and half-kernel
equations hold as equalities of torsion sections. Multiplication
by $d$ recovers $\eta$ everywhere. Both auxiliary $A$-markings
are zero. This gives the claimed universal object and forgetful map.

To check the exact representing clause, first note that two
identifications of a lower object over a point of a reduced scheme
induce the same map on its abelian scheme. On any geometric fiber
their quotient is a polarized $B$-linear automorphism fixing the
lower marking. By PUN6, applied on the completed lattice via the
lower Tate witness, fixing the conjugate-half marking modulo $m$
and preserving the perfect alternating form imply fixing the whole
$m$-torsion. Choose an odd prime $p\mid m$. The automorphism acts
by an element of $1+p\operatorname{End}_{\mathbb Z_p}(T_p A)$.
The polarized automorphism group is finite: in characteristic zero
a polarized complex abelian variety has a faithful integral homology
action preserving the positive definite metric of the polarization
and its complex structure. The integral matrices in this compact
group are finite. The same statement for an individual
characteristic-zero geometric fiber and its proposed automorphism
follows by descending the variety, polarization and endomorphism
to a finitely generated subfield and embedding that field in
$\mathbb C$. The matrix logarithm and exponential are inverse on
this $p$-adic congruence subgroup for odd $p$, so it is torsion free.
Thus the automorphism acts identically on $T_p A$ and on $A$:
over $\mathbb C$ the union of the $p$-power division points is
dense in the torus, and equality descends to the original field.
The two identifications agree over a reduced base because the
source abelian scheme is smooth over that base and is reduced, and
their equalizer is closed. The map on the abelian scheme is all
that is needed to transport the marking.

Now let a reduced test scheme carry an upper object. Multiplication
of its main marking by $d$ restricts it to the lower level, by the
same half-idempotent and tower calculation. The lower representing
law gives a unique map to $X$ and a cartesian identification with
the pulled-back universal family. Transport the $d$-projection of
the upper marking by that identification. By A73H1 it induces
$\alpha$ on every fiber. These are actual morphisms of finite
\'etale sheaves: the marking sections specify the images of
the elements of the finite module, and all module relations may
be checked on geometric fibers. The resulting morphism is an
isomorphism by the same fiber criterion. In particular the native
marking was only required a priori to be additive with the stated
kernel. Its $B$-linearity here follows from the geometric Tate
witness and equality on geometric fibers. Its independence of the
identification follows from the rigidity just proved.
It gives the required unique map to $P$, and (A73.1) identifies
the original marking with the pullback of the universal one.
Finally, two maps from the dual numbers in the infinitesimal clause
of $\mathrm{Moduli}$, one constant, have equal composites to $X$ by
the corresponding lower infinitesimal clause applied to the forgotten
object. They have the same closed point. Formal unramifiedness of
$P\to X$ gives equality of the maps. This proves the last field of
the native moduli record.
\end{proof}
% END A73H2

RTG10 supplies the specified higher coefficient record over $M$.
Base extension of (A73.2) to this field is again finite \'etale
surjective, with the same pullback and marking identities. The
choice of a normalized higher analytic dictionary for these models
is a separate part of the proof of $\mathrm{CoversAt}$.
